\scanpage{5}
de $A$ est fini sur $A$, et d'ailleurs un anneau local. Appliquant Cor. 2 du lemme, on est ramené au cas $A$ \uline{local} \uline{complet} \inserted{\uline{normal}} de dim $\geq\deleted{\unsure{2}}\unsure{3|2}$, donc $X'$ normal de dim $\geq 2$. Alors
\[
 H^0(\widehat X',\underline{\mathcal O}_{X'})\xleftarrow{\sim}
 H^0(X',\underline{\mathcal O}_{X'})\qquad(\text{\unsure{cf|cif|cf.}})
\]
d'où la conclusion.

\uline{Lemme 3} \quad $A$ anneau local \deleted{\unsure{complet}} \inserted{\unsure{noeth.}} complet. \deleted{\unsure{Supposons} $f$ \unsure{appartenant} \unsure{dans les points} \unsure{de $A$|$\in\mathfrak r A$} \unsure{génériques} \illegible{2}.}

$X$, $X'$, \deleted{[$Y'=V(f)\cap X'$]}. Supposons \inserted{$k\geq 1$ et} que

(i) les composantes de $X'$ de dim $\geq k+1$

(ii) $X'$ est connexe en dim $\geq k$

Alors \inserted{$\exists\ \Phi$ fini de $X'$ tel que $f\in\unsure{\mathfrak r A|\mathfrak r(A)}$, $V(f)\cap\Phi=\varnothing$} \deleted{$\exists f\in\mathfrak rA$ tel que} $\Longrightarrow Y'=X'\cap V(f)$ \deleted{Alors $Y'$ est \unsure{de sorte} \unsure{que} \unsure{de dim|en dim} $\geq k$,} \deleted{\unsure{et|est}} connexe en dim $\geq k-1$.

De façon précise, on construit une partie finie $\Phi$ de $X'$ telle que $V(f)\cap\Phi=\varnothing\Longrightarrow$ conclusion voulue.

Il suffit par le cor 1 : lemme 1 que $V(f)\cap\unsure{X'_i|X'}$ soient connexes en dim $\geq k-1$.
% archive sheet number: 4

